\documentclass[aspectratio=169,12pt]{beamer}

% ============================================================
% Theme and visual style
% ============================================================

\usetheme{Madrid}
\usecolortheme{whale}
\usefonttheme{professionalfonts}

\definecolor{MainBlue}{RGB}{45,55,155}
\definecolor{SoftBlue}{RGB}{238,241,250}
\definecolor{DarkText}{RGB}{35,35,40}

\setbeamercolor{structure}{fg=MainBlue}
\setbeamercolor{frametitle}{fg=white,bg=MainBlue}
\setbeamercolor{title}{fg=white}
\setbeamercolor{block title}{fg=white,bg=MainBlue}
\setbeamercolor{block body}{fg=DarkText,bg=SoftBlue}

\setbeamertemplate{navigation symbols}{}

\setbeamertemplate{itemize item}{\small\raise0.2ex\hbox{$\bullet$}}
\setbeamertemplate{itemize subitem}{\scriptsize\raise0.2ex\hbox{$\circ$}}

\setbeamerfont{frametitle}{series=\bfseries}
\setbeamerfont{block title}{series=\bfseries}

% ============================================================
% Packages
% ============================================================

\usepackage[utf8]{inputenc}
\usepackage[T1]{fontenc}
\usepackage[italian]{babel}

\usepackage{amsmath,amssymb,amsfonts}
\usepackage{mathtools}
\usepackage{booktabs}

% ============================================================
% Commands
% ============================================================

\newcommand{\CL}[1]{\llbracket #1 \rrbracket}
\newcommand{\Ag}{\mathcal{A}g}

% ============================================================
% Metadata
% ============================================================

\title{Coalition Logic and Quantification}
\subtitle{From Modal Logic to Quantified Coalition Logic}
\author{Luca Sforza}
\institute{Logic and Reasoning 2025/2026}
\date{\today}

% ============================================================
% Document
% ============================================================

\begin{document}

% ============================================================
% Title
% ============================================================

\begin{frame}
	\titlepage
\end{frame}

\note{
	Good morning everyone. I am Luca Sforza, and today I will present
	my work on Coalition Logic and its quantified extension.

	We will start from classical modal logic, move through the foundations
	of Pauly's Coalition Logic, and finally arrive at Quantified Coalition Logic.

	The goal is to show how logic can be used to model coalitional power
	in multi-agent systems.
}

% ============================================================
% Outline
% ============================================================

\begin{frame}{Outline}
	\begin{enumerate}
		\item Introduction: From Kripke Models to Game Frames
		\item The Concept of Coalitional Power
		\item Coalition Logic (CL): Syntax, Semantics, Axiomatization
		\item The Characterization Theorem
		\item The Succinctness Problem
		\item Quantified Coalition Logic (QCL): Predicates and Operators
		\item Applications: Voting and Social Choice
		\item Conclusions and Reflections
	\end{enumerate}
\end{frame}

% ============================================================
\section{Introduction: From Kripke Models to Game Frames}
% ============================================================

\begin{frame}{Single-Agent Modal Logic}

	\begin{itemize}
		\item
		      In the simplest setting, we consider \textbf{one agent}
		      choosing among different actions.

		\item
		      A \textbf{Kripke model} $(S,R)$ uses an accessibility relation
		      $R$ that associates each state with the states reachable from it.

		\item
		      The modal formula
		      \[
			      \Diamond\varphi
		      \]
		      means that the agent can act in such a way that
		      $\varphi$ becomes true.
	\end{itemize}

	\vspace{0.5em}

	\begin{block}{Semantics}
		\[
			M,s\models\Diamond\varphi
			\iff
			\exists t\in S:
			\; sRt
			\text{ and }
			M,t\models\varphi.
		\]
	\end{block}

\end{frame}

% ============================================================

\begin{frame}{The Multi-Agent Problem}

	\begin{itemize}

		\item
		      Extending the model with individual accessibility relations
		      $R_i$ still considers agents largely \textbf{in isolation}.

		\item
		      In a multi-agent setting, however, agents' actions
		      \textbf{interact}.

		\item
		      What happens when agent $1$ and agent $2$ act
		      \textbf{simultaneously}?

		\item
		      The resulting outcome depends on the
		      \textbf{combination of all participants' choices}.
	\end{itemize}

	\vspace{0.7em}

	\begin{alertblock}{Problem}
		Standard multi-agent modal logic cannot capture this interaction.
	\end{alertblock}

\end{frame}

% ============================================================

\begin{frame}{The Solution: Game Frames (Pauly 2002)}

	\begin{block}{Definition --- Game Frame}
		Marc Pauly introduces \textbf{game frames}:
		at every state, a \textbf{strategic game} is played among the agents,
		and the outcomes of the game are states of the model.
	\end{block}

	\vspace{0.5em}

	\begin{itemize}
		\item
		      Game frames are essentially
		      \textbf{extensive games with simultaneous moves}.

		\item
		      They generalize both Kripke models
		      (the one-player case)
		      and non-normal modal logics
		      (the two-player case).
	\end{itemize}

\end{frame}

% ============================================================

\begin{frame}{Strategic Games}

	Before defining game frames formally, we first introduce
	the strategic games associated with their states.

	\vspace{0.6em}

	A strategic game
	\[
		G=(N,\{\Sigma_i\mid i\in N\},o,S)
	\]
	consists of:

	\begin{description}
		\item[$N$]
		      a finite set of agents (players);

		\item[$\Sigma_i$]
		      the set of strategies available to agent $i$;

		\item[$S$]
		      the set of possible outcome states;

		\item[$o$]
		      an outcome function
		      \[
			      o:\prod_{i\in N}\Sigma_i\longrightarrow S.
		      \]
	\end{description}

	\vspace{0.4em}

	We consider \emph{game forms} only:
	no preference relations are assumed.

\end{frame}

% ============================================================

\begin{frame}{Game Frames}

	\begin{block}{Definition}
		Let $\Gamma^N_S$ denote the set of all strategic games
		over agents $N$ and states $S$.

		A \emph{game frame} is a pair
		\[
			(S,\gamma),
		\]
		where
		\[
			\gamma:S\longrightarrow\Gamma^N_S
		\]
		assigns a strategic game to each state.
	\end{block}

	\vspace{0.5em}

	\begin{itemize}
		\item
		      At every state, the agents play a game.

		\item
		      The outcome leads to a new state,
		      where another game is associated with that state.

		\item
		      $\gamma$ is a \textbf{total function}:
		      every state has an associated game.
	\end{itemize}

\end{frame}

% ============================================================
\section{Coalitional Power}
% ============================================================

\begin{frame}{From Strategy to Power: Effectivity Functions}

	\begin{block}{Definition}
		The \emph{effectivity function}
		\[
			E_G:\mathcal{P}(N)\longrightarrow
			\mathcal{P}(\mathcal{P}(S))
		\]
		derived from a game $G$ is defined by
		\[
			X\in E_G(C)
			\iff
			\exists\sigma_C\,
			\forall\sigma_{N\setminus C}:
			\;
			o(\sigma_C,\sigma_{N\setminus C})\in X.
		\]
	\end{block}

	\begin{itemize}
		\item
		      Coalition $C$ is \textbf{effective for $X$}
		      if it has a joint strategy that forces the outcome into $X$,
		      regardless of what the other agents do.

		\item
		      More generally, for a set of agents $N$ and outcomes $S$,
		      any function
		      \[
			      E:\mathcal{P}(N)\longrightarrow
			      \mathcal{P}(\mathcal{P}(S))
		      \]
		      is called an \emph{effectivity function}.
	\end{itemize}

\end{frame}

% ============================================================

\begin{frame}{Playability (Pauly 2002)}

	Effectivity functions may be required to satisfy different properties,
	depending on the situation we want to model.

	\vspace{0.6em}

	\small

	\begin{description}

		\item[Outcome monotonicity]
		      \[
			      X\in E(C,s),\;
			      X\subseteq X'
			      \;\Rightarrow\;
			      X'\in E(C,s).
		      \]

		\item[Coalition monotonicity]
		      \[
			      C\subseteq C'\subseteq N
			      \;\Rightarrow\;
			      E(C)\subseteq E(C').
		      \]

		\item[$C$-regularity]
		      \[
			      X\in E(C)
			      \;\Rightarrow\;
			      \overline{X}\notin E(\overline{C}).
		      \]

		\item[$C$-maximality]
		      \[
			      \overline{X}\notin E(\overline{C})
			      \;\Rightarrow\;
			      X\in E(C).
		      \]

		\item[Regularity / maximality]
		      $E$ is regular (maximal) if it is
		      $C$-regular ($C$-maximal) for every $C\subseteq N$.

		\item[Superadditivity]
		      If $C_1\cap C_2=\emptyset$, then
		      \[
			      X_1\in E(C_1),\;
			      X_2\in E(C_2)
			      \Rightarrow
			      X_1\cap X_2\in E(C_1\cup C_2).
		      \]

	\end{description}

\end{frame}

% ============================================================

\begin{frame}{Playability (Pauly 2002)}

	The relevant properties depend on the situation we want to model.

	In our setting, the central question is:

	\begin{center}
		\emph{Which effectivity functions arise from some strategic game?}
	\end{center}

	\vspace{0.5em}

	An effectivity function $E$ is \textbf{playable} if:

	\begin{enumerate}
		\item
		      $\forall C\subseteq N,\;
			      \emptyset\notin E(C)$;

		\item
		      $\forall C\subseteq N,\;
			      S\in E(C)$;

		\item
		      $E$ is $N$-maximal;

		\item
		      $E$ is outcome-monotonic;

		\item
		      $E$ is superadditive.
	\end{enumerate}

\end{frame}

% ============================================================

\begin{frame}{Characterization Theorem (Pauly 2002)}

	\vfill

	\begin{theorem}
		An effectivity function $E$ is playable if and only if
		there exists a strategic game $G$ such that
		\[
			E_G=E.
		\]
	\end{theorem}

	\vfill

\end{frame}

% ============================================================

\begin{frame}{From Coalition Frames to Models}

	Let
	\[
		E:
		S
		\longrightarrow
		\mathcal{P}(N)
		\longrightarrow
		\mathcal{P}(\mathcal{P}(S)).
	\]

	\begin{description}

		\item[Coalition Frame]
		      A pair
		      \[
			      \mathcal{F}=(S,E)
		      \]
		      where $E$ assigns a playable effectivity function
		      to each state.

		\item[Coalition Model]
		      A triple
		      \[
			      \mathcal{M}=(S,E,\pi)
		      \]
		      obtained by adding a valuation $\pi$.

	\end{description}

	\vspace{0.3em}

	\[
		\boxed{
			\varphi ::=
			\top
			\mid p
			\mid \neg\varphi
			\mid \varphi\lor\varphi
			\mid \langle\!\langle C\rangle\!\rangle\varphi
			\mid [P]\varphi
		}
	\]

	\vspace{0.3em}

	The semantics of the coalition modality is

	\[
		\mathcal{M},s
		\models
		\langle\!\langle C\rangle\!\rangle\varphi
		\iff
		\varphi^{\mathcal{M}}
		\in E(s,C),
	\]

	where

	\[
		\varphi^{\mathcal{M}}
		=
		\{t\in S\mid\mathcal{M},t\models\varphi\}.
	\]

\end{frame}

\note{
	The characterization theorem ensures that coalition frames correspond
	precisely to game frames.
}

% ============================================================
\section{Coalition Logic}
% ============================================================

\begin{frame}{CL Axiomatization}

	\small

	\textbf{Sound and complete for weak playability models}
	\hfill
	{\scriptsize Ågotnes et al.\ (2008), Table 1}

	\vspace{0.7em}

	\begin{description}

		\item[$\mathbf{Ag\bot}$]
		      \[
			      \neg
			      \langle\!\langle\Ag\rangle\!\rangle
			      \bot
		      \]

		\item[$\mathbf{\top}$]
		      \[
			      \neg
			      \langle\!\langle\emptyset\rangle\!\rangle
			      \bot
			      \to
			      \langle\!\langle C\rangle\!\rangle
			      \top
		      \]

		\item[$\mathbf{\bot}$]
		      \[
			      \langle\!\langle C\rangle\!\rangle\bot
			      \to
			      \langle\!\langle C'\rangle\!\rangle\bot
			      \qquad
			      (C'\subseteq C)
		      \]

		\item[$\mathbf{Ag}$]
		      \[
			      \neg
			      \langle\!\langle\emptyset\rangle\!\rangle
			      \neg\varphi
			      \to
			      \langle\!\langle\Ag\rangle\!\rangle
			      \varphi
		      \]

		\item[$\mathbf{S}$]
		      \[
			      \begin{aligned}
				       &
				      \bigl(
				      \langle\!\langle C_1\rangle\!\rangle\varphi_1
				      \land
				      \langle\!\langle C_2\rangle\!\rangle\varphi_2
				      \bigr)
				      \\
				       & \qquad\to
				      \langle\!\langle C_1\cup C_2\rangle\!\rangle
				      (\varphi_1\land\varphi_2),
				      \qquad
				      C_1\cap C_2=\emptyset.
			      \end{aligned}
		      \]

	\end{description}

	\vspace{0.2em}
	\hrule
	\vspace{0.4em}

	\textbf{Rules:}
	Modus Ponens (MP), all propositional tautologies (Prop), and Distribution:

	\[
		\vdash\varphi\to\psi
		\quad\Rightarrow\quad
		\vdash
		\langle\!\langle C\rangle\!\rangle\varphi
		\to
		\langle\!\langle C\rangle\!\rangle\psi.
	\]

\end{frame}

\note{
	The axiomatization of CL is sound and complete for weak playability models.

	Each axiom encodes one of the semantic playability properties
	at the logical level.

	The rules are Modus Ponens and Distribution.

	The most important axiom for complexity is superadditivity:
	it is precisely what makes satisfiability PSPACE-complete rather than
	NP-complete.

	Without it, each
	$\langle\!\langle C\rangle\!\rangle\varphi$
	requires only one recursive check.
	With it, arbitrary subsets of modal formulas can be combined,
	potentially forcing $2^n$ recursive calls.
}

% ============================================================
\section{Succinctness}
% ============================================================

\begin{frame}{The Succinctness Problem}

	To express

	\begin{center}
		\emph{``some coalition can force $\varphi$''}
	\end{center}

	in Coalition Logic, we need:

	\[
		\bigvee_{C\subseteq\Ag}
		\langle\!\langle C\rangle\!\rangle\varphi.
	\]

	\vspace{0.8em}

	\begin{itemize}

		\item
		      With $n$ agents, the formula contains
		      \textbf{$2^n$ disjuncts}.

		\item
		      Its size therefore grows
		      \textbf{exponentially} with the number of agents.

		\item
		      General properties such as
		      \emph{``a majority can win''}
		      become practically impossible to express compactly in CL.

	\end{itemize}

\end{frame}

% ============================================================
\section{Quantified Coalition Logic}
% ============================================================

\begin{frame}{The Innovation: Predicates over Coalitions}

	Instead of explicitly enumerating coalition members,
	QCL describes coalitions using \textbf{predicates} $P$:

	\[
		P ::=
		\operatorname{subseteq}(C)
		\mid
		\operatorname{supseteq}(C)
		\mid
		\neg P
		\mid
		P\lor P.
	\]

	\vspace{0.5em}

	\begin{align*}
		C\models\operatorname{subseteq}(D)
		 & \iff C\subseteq D,
		\\[0.3em]
		C\models\operatorname{supseteq}(D)
		 & \iff C\supseteq D.
	\end{align*}

	\vspace{0.4em}

	\textbf{Derived predicates:}

	\begin{itemize}
		\item
		      $\operatorname{eq}(D)
			      \equiv
			      \operatorname{subseteq}(D)
			      \land
			      \operatorname{supseteq}(D)$
		      \hfill exact equality;

		\item
		      $\operatorname{incl}(i)
			      \equiv
			      \operatorname{supseteq}(\{i\})$
		      \hfill agent $i$ is a member;

		\item
		      $\operatorname{excl}(i)
			      \equiv
			      \neg\operatorname{incl}(i)$
		      \hfill agent $i$ is not a member;

		\item
		      $\operatorname{any}
			      \equiv
			      \operatorname{supseteq}(\emptyset)$
		      \hfill any coalition.
	\end{itemize}

\end{frame}

\note{
	Quantified Coalition Logic addresses the succinctness problem
	by introducing a quantification mechanism over coalitions.

	Instead of enumerating the members of a coalition,
	coalitions are described through predicates.

	The two atomic predicates are inclusion and containment:
	subseteq(D) checks whether the coalition is a subset of D,
	while supseteq(D) checks whether it contains D.

	From these primitives we derive useful predicates:
	eq(D) for exact equality, incl(i) to test whether agent i belongs
	to the coalition, excl(i) for the opposite condition,
	and any, which is satisfied by every coalition.

	subseteq and supseteq are mutually definable because the set of agents
	is finite. However, defining supseteq starting from subseteq requires
	an exponentially long formula, so both are needed as primitives.
}

% ============================================================

\begin{frame}{Syntax of QCL}

	\vfill

	\[
		\boxed{
			\varphi ::=
			\top
			\mid p
			\mid \neg\varphi
			\mid \varphi\lor\varphi
			\mid \langle P\rangle\varphi
			\mid [P]\varphi
		}
	\]

	\vspace{1.2em}

	\begin{itemize}

		\item
		      $\langle P\rangle\varphi$:
		      \textbf{some} coalition satisfying $P$
		      can force $\varphi$.

		      \vspace{0.5em}

		\item
		      $[P]\varphi$:
		      \textbf{every} coalition satisfying $P$
		      can force $\varphi$.

	\end{itemize}

	\vfill

\end{frame}

% ============================================================

\begin{frame}{Semantics of QCL}

	\small

	\begin{block}{Propositional cases}

		\[
			\begin{aligned}
				\mathcal{M},s\models_{\mathrm{QCL}}\top
				 &      \\[0.25em]
				\mathcal{M},s\models_{\mathrm{QCL}}p
				 & \iff
				p\in\pi(s)
				\\[0.25em]
				\mathcal{M},s\models_{\mathrm{QCL}}\neg\varphi
				 & \iff
				\mathcal{M},s\not\models_{\mathrm{QCL}}\varphi
				\\[0.25em]
				\mathcal{M},s\models_{\mathrm{QCL}}\varphi\lor\psi
				 & \iff
				\mathcal{M},s\models_{\mathrm{QCL}}\varphi
				\text{ or }
				\mathcal{M},s\models_{\mathrm{QCL}}\psi
			\end{aligned}
		\]

	\end{block}

	\begin{block}{Coalition modalities}

		\[
			\begin{aligned}
				\mathcal{M},s
				\models_{\mathrm{QCL}}
				\langle P\rangle\varphi
				\iff {} &
				\exists C\subseteq\Ag:
				C\models_{\mathrm{cp}}P
				\land
				\exists S\in E(C,s):
				\\
				        &
				\qquad
				\forall s'\in S,\;
				\mathcal{M},s'
				\models_{\mathrm{QCL}}
				\varphi
			\end{aligned}
		\]

		\[
			\begin{aligned}
				\mathcal{M},s
				\models_{\mathrm{QCL}}
				[P]\varphi
				\iff {} &
				\forall C\subseteq\Ag:
				C\models_{\mathrm{cp}}P
				\Longrightarrow
				\exists S\in E(C,s):
				\\
				        &
				\qquad
				\forall s'\in S,\;
				\mathcal{M},s'
				\models_{\mathrm{QCL}}
				\varphi
			\end{aligned}
		\]

	\end{block}

\end{frame}

% ============================================================

\end{document}
